In the absence of a heuristic approach, the search for an optimal augmentation policy is highly empirical and dependent on the specific dataset. The primary objective of this study is to evaluate suitable augmentation techniques for the marine wildlife dataset by employing basic augmentation methods and exploring the effectiveness of GAN-generated images to improve the classification of underrepresented species.  Additionally, this study aims to investigate whether an augmentation technique that enhances the performance of one task might necessarily be beneficial for another, and whether the effectiveness of the augmentation policy is model-agnostic. To sum up, this study seeks to address the following questions:

\begin{itemize}
\item How is the effectiveness of common augmentation (geometric, color-space, cutout) for the maritime wildlife aerial image data-set in detection and classification tasks?
\item Is the augmentation effectiveness only dependent by the dataset and agnostic to the model used and the tasks?
\item If and how low-data classes may benefit from the GAN-synthetic method?
\end{itemize}
